\documentclass[10pt,a4paper]{amsart}
\usepackage{amsmath,amssymb,amsthm,graphicx}
\usepackage{verbatim}
\newenvironment{reference}{\comment}{\endcomment}
%\newenvironment{reference}{}{}
\newenvironment{slogan}{\comment}{\endcomment}
\newenvironment{history}{\comment}{\endcomment}

% For commutative diagrams we use Xy-pic
\usepackage[all]{xy}

% We use 2cell for 2-commutative diagrams.
\xyoption{2cell}
% Theorem environments.
%
\theoremstyle{plain}
\newtheorem{theorem}[subsection]{Theorem}
\newtheorem{proposition}[subsection]{Proposition}
\newtheorem{lemma}[subsection]{Lemma}

\theoremstyle{definition}
\newtheorem{definition}[subsection]{Definition}
\newtheorem{example}[subsection]{Example}
\newtheorem{exercise}[subsection]{Exercise}
\newtheorem{situation}[subsection]{Situation}

\theoremstyle{remark}
\newtheorem{remark}[subsection]{Remark}
\newtheorem{remarks}[subsection]{Remarks}

\numberwithin{equation}{subsection}

% Macros
%
\def\lim{\mathop{\mathrm{lim}}\nolimits}
\def\colim{\mathop{\mathrm{colim}}\nolimits}
\def\Spec{\mathop{\mathrm{Spec}}}
\def\Hom{\mathop{\mathrm{Hom}}\nolimits}
\def\Ext{\mathop{\mathrm{Ext}}\nolimits}
\def\SheafHom{\mathop{\mathcal{H}\!\mathit{om}}\nolimits}
\def\SheafExt{\mathop{\mathcal{E}\!\mathit{xt}}\nolimits}
\def\Sch{\mathit{Sch}}
\def\Mor{\mathop{\mathrm{Mor}}\nolimits}
\def\Ob{\mathop{\mathrm{Ob}}\nolimits}
\def\Sh{\mathop{\mathit{Sh}}\nolimits}
\def\NL{\mathop{N\!L}\nolimits}
\def\CH{\mathop{\mathrm{CH}}\nolimits}
\def\proetale{{pro\text{-}\acute{e}tale}}
\def\etale{{\acute{e}tale}}
\def\QCoh{\mathit{QCoh}}
\def\Ker{\mathop{\mathrm{Ker}}}
\def\Im{\mathop{\mathrm{Im}}}
\def\Coker{\mathop{\mathrm{Coker}}}
\def\Coim{\mathop{\mathrm{Coim}}}

% Boxtimes
%
\DeclareMathSymbol{\boxtimes}{\mathbin}{AMSa}{"02}

%
% Macros for moduli stacks/spaces
%
\def\QCohstack{\mathcal{QC}\!\mathit{oh}}
\def\Cohstack{\mathcal{C}\!\mathit{oh}}
\def\Spacesstack{\mathcal{S}\!\mathit{paces}}
\def\Quotfunctor{\mathrm{Quot}}
\def\Hilbfunctor{\mathrm{Hilb}}
\def\Curvesstack{\mathcal{C}\!\mathit{urves}}
\def\Polarizedstack{\mathcal{P}\!\mathit{olarized}}
\def\Complexesstack{\mathcal{C}\!\mathit{omplexes}}
% \Pic is the operator that assigns to X its picard group, usage \Pic(X)
% \Picardstack_{X/B} denotes the Picard stack of X over B
% \Picardfunctor_{X/B} denotes the Picard functor of X over B
\def\Pic{\mathop{\mathrm{Pic}}\nolimits}
\def\Picardstack{\mathcal{P}\!\mathit{ic}}
\def\Picardfunctor{\mathrm{Pic}}
\def\Deformationcategory{\mathcal{D}\!\mathit{ef}}
\usepackage{cleveref}
\crefname{originallemma}{original lemma}{original lemmas}
\Crefname{originallemma}{Original lemma}{Original lemmas}
\makeatletter
\newlabel{formal-spaces-lemma-formal-completion-types}{{0AQ1}{source}{}{sourcepointer-1}{}}
\newlabel{formal-spaces-lemma-formal-completion-types@cref}{{[originallemma][0][]0AQ1}{[1][0][]source}{}{}{}}
\makeatother
\crefname{originalreference}{original reference}{original references}
\Crefname{originalreference}{Original reference}{Original references}
\makeatletter
\newlabel{lemma-completion-proper-surjective-rig-surjective}{{0AQW}{source}{}{sourcepointer-2}{}}
\newlabel{lemma-completion-proper-surjective-rig-surjective@cref}{{[originalreference][0][]0AQW}{[1][0][]source}{}{}{}}
\makeatother
\crefname{originalreference}{original reference}{original references}
\Crefname{originalreference}{Original reference}{Original references}
\makeatletter
\newlabel{lemma-etale-gives-rig-etale}{{0AR2}{source}{}{sourcepointer-3}{}}
\newlabel{lemma-etale-gives-rig-etale@cref}{{[originalreference][0][]0AR2}{[1][0][]source}{}{}{}}
\makeatother
\crefname{originalreference}{original reference}{original references}
\Crefname{originalreference}{Original reference}{Original references}
\makeatletter
\newlabel{lemma-proper-gives-proper}{{0GDD}{source}{}{sourcepointer-4}{}}
\newlabel{lemma-proper-gives-proper@cref}{{[originalreference][0][]0GDD}{[1][0][]source}{}{}{}}
\makeatother
\crefname{originalreference}{original reference}{original references}
\Crefname{originalreference}{Original reference}{Original references}
\makeatletter
\newlabel{lemma-separated-mono-open-diagonal-rig-surjective}{{0GDI}{source}{}{sourcepointer-5}{}}
\newlabel{lemma-separated-mono-open-diagonal-rig-surjective@cref}{{[originalreference][0][]0GDI}{[1][0][]source}{}{}{}}
\makeatother
\crefname{originalsection}{original section}{original sections}
\Crefname{originalsection}{Original section}{Original sections}
\makeatletter
\newlabel{section-completion-functor}{{0GDB}{source}{}{sourcepointer-6}{}}
\newlabel{section-completion-functor@cref}{{[originalsection][0][]0GDB}{[1][0][]source}{}{}{}}
\makeatother
\title{Proper modifications algebraize to formal modifications}
\author{The Stacks Project Authors}
\date{stacks-project:restricted.tex:local-snapshot-20261012; GNU FDL 1.2 or later}
\begin{document}
\maketitle
\noindent\textit{Editorial note.} All delivered mathematical text is the author's own, copied from the pinned Stacks Project chapter file restricted.tex; no statement or proof step was rewritten, reordered or removed.\par\medskip
\noindent\textit{Editorial note.} The counted proof is the author's complete argument: the counted statement is the author lemma environment labelled lemma-modification-gives-formal-modification; the proof environment immediately following it is its own proof\par\medskip
\noindent\textit{Editorial note.} The label lemma-modification-gives-formal-modification names the counted statement, and the proof environment immediately following it in the source is delivered verbatim as its proof.\par\medskip
\noindent\textit{Editorial note.} The counted statement is carried by the author's own lemma environment in the chapter source.\par\medskip
\noindent\textit{Allowed standard prerequisites (editorial).}\begin{itemize}
\item the definition of a formal modification f\_\{/T'\} : X'\_\{/T'\} to X\_\{/T\} via the four conditions: locally Noetherian source and target, properness, rig-etale-ness, and rig-surjectivity including on diagonals (section-completion-functor setup)
\item the four lemmas of the same chapter: proper-gives-proper, etale-gives-rig-etale, completion-proper-surjective-rig-surjective, and separated-mono-open-diagonal-rig-surjective
\item the Formal Spaces lemma that formal completions of locally Noetherian schemes are locally Noetherian formal algebraic spaces (formal-spaces-lemma-formal-completion-types)
\end{itemize}\medskip
\ifdefined\hypertarget\hypertarget{sourcepointer-1}{}\fi
\noindent\textit{Source reference.} Original lemma 0AQ1: sources/stacks-chow/tags.txt, source lines 13494--13494. Prerequisite pointer invoked at proof line(s) 6355; it names a standard fact of a sibling chapter that the counted argument uses, and every step of the counted proof built on it is retained verbatim in the delivered body.\par\medskip
\ifdefined\hypertarget\hypertarget{sourcepointer-2}{}\fi
\noindent\textit{Source reference.} Original reference 0AQW: sources/stacks-restricted/pinned-original-source.tex, source lines 6201--6201. Prerequisite pointer invoked at proof line(s) 6361; it names a standard fact of the same chapter that the counted argument uses, and every step of the counted proof built on it is retained verbatim in the delivered body.\par\medskip
\ifdefined\hypertarget\hypertarget{sourcepointer-3}{}\fi
\noindent\textit{Source reference.} Original reference 0AR2: sources/stacks-restricted/pinned-original-source.tex, source lines 6186--6186. Prerequisite pointer invoked at proof line(s) 6360; it names a standard fact of the same chapter that the counted argument uses, and every step of the counted proof built on it is retained verbatim in the delivered body.\par\medskip
\ifdefined\hypertarget\hypertarget{sourcepointer-4}{}\fi
\noindent\textit{Source reference.} Original reference 0GDD: sources/stacks-restricted/pinned-original-source.tex, source lines 6053--6053. Prerequisite pointer invoked at proof line(s) 6359; it names a standard fact of the same chapter that the counted argument uses, and every step of the counted proof built on it is retained verbatim in the delivered body.\par\medskip
\ifdefined\hypertarget\hypertarget{sourcepointer-5}{}\fi
\noindent\textit{Source reference.} Original reference 0GDI: sources/stacks-restricted/pinned-original-source.tex, source lines 6241--6241. Prerequisite pointer invoked at proof line(s) 6362; it names a standard fact of the same chapter that the counted argument uses, and every step of the counted proof built on it is retained verbatim in the delivered body.\par\medskip
\ifdefined\hypertarget\hypertarget{sourcepointer-6}{}\fi
\noindent\textit{Source reference.} Original section 0GDB: sources/stacks-restricted/pinned-original-source.tex, source lines 5953--5953. Prerequisite pointer invoked at proof line(s) 6349; it names a standard fact of the same chapter that the counted argument uses, and every step of the counted proof built on it is retained verbatim in the delivered body.\par\medskip
\begin{lemma}
\label{lemma-modification-gives-formal-modification}
Let $S$, $f : X' \to X$, $T \subset |X|$, $U \subset X$,
$T' \subset |X'|$, and $U' \subset X'$ be as in
Section \ref{section-completion-functor}.
If $X$ is locally Noetherian, $f$ is proper, and $U' \to U$ is an isomorphism,
then $f_{/T} : X'_{/T'} \to X_{/T}$ is a formal modification.
\end{lemma}

\begin{proof}
By Formal Spaces, Lemmas \ref{formal-spaces-lemma-formal-completion-types}
the source and target of the arrow are locally Noetherian
formal algebraic spaces.
The other conditions follow from
Lemmas \ref{lemma-proper-gives-proper},
\ref{lemma-etale-gives-rig-etale},
\ref{lemma-completion-proper-surjective-rig-surjective}, and
\ref{lemma-separated-mono-open-diagonal-rig-surjective}.
\end{proof}

\end{document}
